\paragraph{Feasible actions versus learnable utility.}
Figure~\ref{fig:capacity-discrimination} separates the candidate menu,
chronological credit and causal discrimination. All four hindsight
arms use the same total quota and original risk. Chronological credits
reduce feasible positive plans, but the strict menu still permits
\DirectCountStrict{}/\DirectCountViews{} verified plans. A learned
online count policy reaches \DirectCountCount{} under the same strict
constraints. The gap is not removed by training on token-feasible plans
(supplement). We then learn direct gains and risks for each nested
action, keeping nominations, architecture and online budget fixed.
The motion-only and peer models have gain AUROCs \DirectCountMotionAUC{}
and \DirectCountPeerAUC{} and joint pass counts \DirectCountMotion{}
and \DirectCountPeer{}, against \DirectCountProxy{} for a simple
observed-prefix score. At exactly \DirectCountMatchedEdits{} edits in a
retrospective match, the peer effect is \DirectCountGap{} pp
[\DirectCountLow{}, \DirectCountHigh{}], and joint passes fall from
\DirectCountMatchedMotion{} to \DirectCountMatchedPeer{}. The evidence
does not show that the added context recovers the feasible capacity.

\begin{figure*}[t]
\centering
\includegraphics[width=.9\textwidth]{capacity_discrimination.pdf}
\caption{OpenDD development. (a) Four label-aware feasible plan counts
and an online count controller. Full-menu results are shared across the
three seeds; 606 views represent 202 distinct recordings. (b) The peer
head compared with motion and an observed-prefix control. Error bars
resample locations after averaging seeds. Exact realized count matching
uses all histories in a recording but no labels, and is not online.}
\label{fig:capacity-discrimination}
\end{figure*}
